%==========================================================================
% CAPÍTULO 6 — CONCLUSÃO E TRABALHOS FUTUROS
%==========================================================================
\chapter{Conclusão e Trabalhos Futuros}
\label{chap:conclusao}

\section{Conclusão}

O presente estágio teve como objetivo investigar a viabilidade de um sistema
de \textbf{deteção automática de gomos em videiras} utilizando técnicas de
visão computacional e \textit{deep learning}, com vista à sua integração num
futuro sistema robótico de poda mecanizada.

Ao longo do período de estágio foram desenvolvidas duas abordagens
complementares. O \textbf{Modelo~1}, baseado num dataset espanhol de domínio
público, permitiu validar a \textit{pipeline} técnica e demonstrar que o
    YOLO26 Nano é capaz de detetar gomos em condições de imagem moderadamente
controladas. O \textbf{Modelo~2}, construído a partir de vídeos captados nas
vinhas de Palmela, representa um esforço mais ambicioso e realista, visando
treinar um modelo robusto em condições próximas de uma operação robótica real.

Os principais resultados e contribuições do estágio são:

\begin{itemize}
  \item Construção de uma metodologia de aquisição e anotação de imagens
    adaptada à complexidade do ambiente vitícola português;

  \item Desenvolvimento de um dataset original com aproximadamente
    \textbf{400 imagens anotadas} de vinhas da região de Palmela, com
    quatro classes:     \textit{gomo}, \textit{sarmento}, \textit{tronco} e
    \textit{objeto};

  \item Identificação e documentação dos principais \textbf{desafios
    técnicos} associados à deteção de gomos em cenas reais: oclusões, fundo
    denso, variações de iluminação e alta densidade de instâncias;

  \item Avaliação da adequação do \textbf{YOLO26 Nano} para esta tarefa,
    com identificação dos seus pontos fortes (eficiência computacional,
    facilidade de integração) e limitações (sensibilidade ao volume e
    qualidade dos dados de treino);

  \item Prova de conceito de que, com um dataset suficientemente grande e
    diverso, um modelo de deteção de objetos pode localizar gomos e
    sarmentos em imagens de campo com complexidade real.
\end{itemize}

É importante reconhecer que a obtenção de um modelo plenamente funcional e
robusto para uso em sistemas robóticos reais \textbf{excede o âmbito de um
único semestre de estágio}. A complexidade visual das imagens de Palmela,
combinada com a necessidade de volumes de dados elevados, exige um esforço
continuado. Contudo, as bases estabelecidas neste trabalho são um ponto de
partida sólido e documentado para projetos de continuação.

\section{Trabalhos Futuros}

Com base na experiência adquirida e nas limitações identificadas, propõem-se
as seguintes linhas de trabalho futuro:

\subsection{Expansão e Diversificação do Dataset}

\begin{itemize}
  \item Aumentar o dataset de Palmela para um mínimo de \textbf{1000--2000
    imagens anotadas}, cobrindo diferentes épocas do ano (pré-poda, pós-
    poda, início do abrolhamento) e diferentes condições meteorológicas.

  \item Incorporar imagens de outras regiões vitivinícolas portuguesas e
    de diferentes cultivares, para aumentar a diversidade do modelo.

  \item Explorar a \textbf{anotação semi-automática}: usar um modelo
    preliminar para gerar sugestões de anotação que sejam depois revistas
    manualmente, reduzindo o esforço por imagem.
\end{itemize}

\subsection{Arquiteturas de Modelo Mais Avançadas}

\begin{itemize}
  \item Avaliar variantes maiores do YOLO26 (Small, Medium) ou modelos da
    geração seguinte (RT-DETR, YOLO-NAS) para determinar se o ganho
    em desempenho justifica o custo computacional adicional.

  \item Explorar a \textbf{segmentação de instâncias} (YOLO26-Seg) em vez
    de apenas \textit{bounding boxes}, o que poderia fornecer informação
    mais precisa sobre a geometria dos sarmentos, útil para o cálculo do
    ponto de corte.

  \item Investigar modelos baseados em \textit{transformers} para visão
    (DETR, Grounding DINO) que possam lidar melhor com oclusões severas.
\end{itemize}

\subsection{Integração com Informação de Profundidade}

A adição de uma câmera de profundidade (RGB-D) ao sistema permitiria:

\begin{itemize}
  \item Separar objetos em primeiro e segundo plano com base na distância,
    reduzindo a confusão entre sarmentos de diferentes planos;

  \item Calcular as coordenadas 3D dos gomos para guiar o braço robótico
    ao ponto de corte preciso.
\end{itemize}

\subsection{Validação em Ambiente Real}

\begin{itemize}
  \item Integrar o modelo desenvolvido num protótipo robótico e realizar
    testes \textit{in situ} nas vinhas de Palmela;

  \item Definir métricas de desempenho operacional (taxa de gomos
    corretamente identificados por metro linear de vinha, tempo de
    processamento por frame).
\end{itemize}

\subsection{Técnicas de \textit{Data Augmentation} Específicas}

\begin{itemize}
  \item Explorar augmentações mais agressivas e específicas ao domínio,
    como simulação de diferentes condições de iluminação (\textit{random
    brightness/contrast}), adição de ruído de movimento (\textit{motion
    blur}) e recortes de imagens densas.

  \item Investigar técnicas de \textit{mosaic augmentation} e \textit{copy-
    paste} para aumentar artificialmente a densidade de instâncias por
    imagem.
\end{itemize}
